\section{Introduction}



% \subsection{Setup}

%   Let $\kkk$ be an algebraically closed field of arbitrary characteristic.
%   Let $X$ be a projective variety of dimension $d$ over $\kkk$, and $f: X \dashrightarrow X$ be a dominant rational map.


% \subsection{Results}

  % Fix a projective (geometrically integral) variety $X$ of dimension $d$ and a dominant rational map $f: X \dashrightarrow X$ over a arbitrary field $\kk$.

  We work over an arbitrary field $\kk$ without any restriction on the characteristic.
  A variety is always assumed to be geometrically integral.

  \begin{mainthm}[{= \cref{thm:big_inequality_by_mu_i}, ref. \cite[Theorem 3.7]{Xie24recursive_inequalities}}]\label{mainthm:recursive_inequality}
    Let $X$ be a projective variety of dimension $d$ and $f: X \dashrightarrow X$ be a dominant rational map.
    For $i = 0, 1, \ldots, d$ and $\varepsilon \in (0, 1)$, there exists $m_\varepsilon > 0$ such that for all $m \geq m_\varepsilon$,
    \[  L_{2m} - \varepsilon^m \mu_i^m L_m + \mu_i^m \mu_{i+1}^m L \]
    is big.
  \end{mainthm}

  \begin{mainthm}[{= \cref{thm:semi-continuity-of-dynamical-degrees}, ref. \cite[Theorem 1.8]{Xie24recursive_inequalities}}]\label{mainthm:semi-continuity}
    Let $S$ be an integral noetherian scheme, and $\pi: X \to S$ be a flat projective morphism. 
    Let $f: X \dashrightarrow X$ be a dominant rational map over $S$.
    Then for every $i$, the function $s \mapsto \lambda_i(f_s)$ is lower semi-continuous.
  \end{mainthm}

  \begin{mainthm}[{= \cref{thm:density_of_periodic_points_for_cohomologically_hyperbolic_maps}, ref. \cite[Theorem 1.12]{Xie24recursive_inequalities}}]\label{mainthm:periodic_points}
    Let $X$ be a projective variety and $f: X \dashrightarrow X$ be a dominant rational map.
    If $f$ is cohomologically hyperbolic, then the set of periodic points of $f$ is Zariski dense in $X$.
  \end{mainthm}

  \begin{mainthm}[{= \Yang{}, ref. \cite[Theorem 1.14]{Xie24recursive_inequalities}}]\label{mainthm:application_to_KSC}
    Let $X$ be a projective surface over $\overline{\bbQ}$ and $f: X \dashrightarrow X$ be a dominant rational map.
    Assume that $\mu_2(f) < 1$ or $\mu_1(f) = \mu_2(f)$.
    Then the Kawaguchi-Silverman conjecture holds for $(X, f)$.
  \end{mainthm}